\documentclass[10pt,letterpaper]{article}
\usepackage{cornell-notes}

\title{Clause 17.11.02.03: Class Weak Ordering}
\author{}
\date{}
\setDocTitle{Clause 17.11.02.03: Class Weak Ordering}
\setDocSubtitle{C++ 2024}
\setDocOwner{C++ 2024 Cornell Notes}
\setDocDate{\today}
\setCornellCollection{C++ 2024 Cornell Notes}
\setCornellUnitType{Clause}
\setCornellUnitNumber{17.11.02.03}
\setCornellUnitTitle{Class Weak Ordering}
\hypersetup{pdftitle={Clause 17.11.02.03: Class Weak Ordering},pdfsubject={Cornell notes},pdfauthor={}}

\begin{document}
\maketitle
\begin{CornellOverviewBox}{Overview}
\textbf{Classification:} Standard library - Normative\\[2pt]
\textbf{Learning objective:} Understand the rules, terminology, and semantic consequences associated with Class weak\_ordering.
\end{CornellOverviewBox}\begin{CornellSummaryBox}{Summary}
{\sffamily\bfseries\color{Accent} Summary}\par\vspace{3pt}Section 17.11.2.3 focuses on Class weak\_ordering. Apply the specified interface together with its mandates, effects, result conditions, exception behavior, and complexity constraints. When applying it, preserve the distinction between mandatory requirements, recommendations, permissions, and behavior deliberately left undefined, unspecified, or implementation-defined.
\end{CornellSummaryBox}

\end{document}
